\documentclass[11pt]{article}
\usepackage[T1]{fontenc}
\usepackage{lmodern}
\usepackage[margin=1in]{geometry}
\usepackage{amsmath,amssymb,amsthm,mathtools}
\usepackage{booktabs,array}
\usepackage{microtype}
\usepackage{xcolor}
\usepackage[colorlinks=true,linkcolor=blue!55!black,citecolor=blue!55!black,urlcolor=blue!55!black]{hyperref}
\hypersetup{pdftitle={The Curveball chain on regular graphs: spectral gap at most 2/n and a comparison with random transpositions},pdfauthor={Mark T. Short}}
\numberwithin{equation}{section}
\newtheorem{theorem}{Theorem}[section]
\newtheorem{lemma}[theorem]{Lemma}
\newtheorem{proposition}[theorem]{Proposition}
\newtheorem{corollary}[theorem]{Corollary}
\newtheorem{conjecture}[theorem]{Conjecture}
\theoremstyle{definition}
\newtheorem{definition}[theorem]{Definition}
\theoremstyle{remark}
\newtheorem{remark}[theorem]{Remark}
\newcommand{\NMAX}{11}
\title{The Curveball chain on regular graphs:\\spectral gap at most $2/n$ and a
comparison with random transpositions}
\author{Mark T. Short\\[2pt]\small Independent researcher\\\small\url{https://seveneves.ai}}
\date{November 13, 2026}
\begin{document}
\maketitle
\begin{abstract}
For the Curveball chain that performs a uniform trade at a uniformly chosen
vertex pair of a $d$-regular graph on $n$ vertices, we prove that $1-2/n$ is
an eigenvalue for every $n$ and $d$, so the spectral gap is at most $2/n$ and
the relaxation time is at least $n/2$ trades. We prove that the chain
dominates one half of the random-transposition walk on vertex labels, which
bounds every eigenvalue outside two isotypic components and shows that the
gap equals $2/n$ unless a smaller eigenvalue has an eigenfunction depending
only on the rooted isomorphism class of the graph. The gap equals $2/n$ for
$d\in\{1,n-2\}$ and, by an exact computer-assisted computation, for every
regular degree sequence with $n\le\NMAX$; equality for all $n$ and $d$ is
open.
\end{abstract}
\section{Introduction}\label{sec:introduction}
Let $n\ge4$ and $1\le d\le n-2$ with $nd$ even. Write $\Omega=\Omega_{n,d}$ for
the set of simple $d$-regular graphs on the vertex set $[n]=\{1,\ldots,n\}$, put
$M=|\Omega|$, and let $\pi$ be the uniform probability measure on $\Omega$. A
\emph{Curveball trade} at a vertex pair $\{x,w\}$ redistributes the vertices
adjacent to exactly one of $x$ and $w$ between these two endpoints, keeping
both endpoint degrees and every other adjacency~\cite{SNBFS14,CBS18}.

We use the following precise discrete-time convention, which is that
of~\cite{S26}. The Curveball chain $K$ chooses an unordered vertex pair
uniformly and moves to a graph chosen uniformly among those obtainable from
the current one by a Curveball trade at that pair, the current graph included.
No additional holding probability is imposed, and time counts trades. For an
unordered pair $a$, let $E_a$ average a function over the graphs obtainable by
a trade at $a$, and put
\begin{equation}\label{eq:operators}
 h_a=I-E_a,\qquad H=\sum_{a\in\binom{[n]}2}h_a,\qquad
 K=I-\frac{H}{\binom n2}.
\end{equation}
The kernel $K$ is symmetric, and $\operatorname{gap}(K)$ is one minus its
second largest eigenvalue; $\operatorname{gap}(H)$ is the least eigenvalue of
$H$ on the orthogonal complement $\mathbf1^\perp$ of the constants in
$L^2(\pi)$. Thus $\operatorname{gap}(K)=\operatorname{gap}(H)/\binom n2$. The
symmetric group $S_n$ acts on $\Omega$ by relabeling vertices, and on functions
by $(T_\sigma f)(G)=f(\sigma^{-1}G)$. Let
\[
 L_T=\sum_{a\in\binom{[n]}2}(I-T_a)
\]
be the Laplacian of the random-transposition walk on $\Omega$, where $T_a$ is
the relabeling by the transposition of the two vertices of $a$. For a
partition $\lambda$ of $n$, let $U_\lambda$ be the $\lambda$-isotypic component
of $L^2(\pi)$ and let
$\operatorname{cont}(\lambda)=\sum_i\bigl[\binom{\lambda_i}2-(i-1)\lambda_i\bigr]$
be the sum of the contents of its cells. A function on $\Omega$ is a
\emph{rooted-class function} if it is invariant under every relabeling that
fixes the vertex $n$; such a function depends only on the isomorphism class of
the graph rooted at $n$. Here and throughout, logarithms are natural.

\begin{theorem}[Regular graphs]\label{thm:main}
Let $n\ge4$ and $1\le d\le n-2$ with $nd$ even.
\begin{enumerate}
\item For every pair $\{x,y\}$, the centered edge indicator
$g_{xy}=\mathbf1[xy]-d/(n-1)$ is a nonzero function with $Hg_{xy}=(n-1)g_{xy}$.
Consequently
\begin{equation}\label{eq:upper}
 \operatorname{gap}(H)\le n-1,\qquad \operatorname{gap}(K)\le\frac2n .
\end{equation}
\item $H\succeq\tfrac12L_T$. Consequently, for every partition $\lambda$ of
$n$,
\begin{equation}\label{eq:sector}
 H\succeq\frac12\Bigl[\binom n2-\operatorname{cont}(\lambda)\Bigr]
 \quad\text{on }U_\lambda .
\end{equation}
The right side is at least $n-1$ for every $\lambda$ other than $(n)$ and
$(n-1,1)$, and equals $n/2$ for $\lambda=(n-1,1)$.
\item Let $\mu$ be the least eigenvalue of $H$ on the centered rooted-class
functions, with $\mu=\infty$ when every rooted-class function is constant.
Then
\begin{equation}\label{eq:reduction}
 \operatorname{gap}(H)=\min\{n-1,\mu\}.
\end{equation}
\item If $d\in\{1,n-2\}$, then $\operatorname{gap}(H)=n-1$ and
$\operatorname{gap}(K)=2/n$.
\end{enumerate}
\end{theorem}

Part~(2) is stated for regular degree sequences. For other sequences the
transposition of two vertices of different degrees does not preserve
$\Omega_d$, and Remark~\ref{cmp:general} gives the comparison for the pairs of
equal degree. No lower bound for $\mu$ beyond the universal bound
$H\succeq2(I-\Pi)$ of~\cite[Theorem~1.1 and Corollary~6.4]{S26} is proved
here for general $n$ and $d$. For $d=1$ the chain is a lazy version of the
walk of Diaconis and Holmes on perfect matchings,
$K=(1-\beta)I+\beta K_{\mathrm{DH}}$ with $\beta=(n-2)/(2(n-1))$, so
part~(4) for $d=1$ follows from their
spectrum~\cite[Theorem~2 and Corollary~1]{DH02}
(Remark~\ref{root:matchings}); for $d=n-2$ it follows from the case $d=1$ by
complementation.

\begin{theorem}[Small cases]\label{thm:small}
Let $1\le d\le n-2$ with $nd$ even. If $4\le n\le\NMAX$, or if $n=12$ and
$d\notin\{5,6\}$, then $\operatorname{gap}(H)=n-1$ and
$\operatorname{gap}(K)=2/n$ on $\Omega_{n,d}$. The lower bound
uses the exact integer computation of Section~\ref{cert:section}.
\end{theorem}

\begin{corollary}[Relaxation and mixing]\label{cor:mixing}
Let $d(s)=\max_{G\in\Omega}\frac12\sum_{G'\in\Omega}|K^s(G,G')-\pi(G')|$ and
$t_{\mathrm{mix}}(\varepsilon)=\min\{s\ge0:d(s)\le\varepsilon\}$. Under the
hypotheses of Theorem~\ref{thm:main}, the relaxation time
$1/\operatorname{gap}(K)$ is at least $n/2$, and for $0<\varepsilon<1/2$,
\begin{equation}\label{eq:mixlower}
 t_{\mathrm{mix}}(\varepsilon)\ge\Bigl(\frac n2-1\Bigr)\log\frac1{2\varepsilon}.
\end{equation}
Whenever $\operatorname{gap}(K)=2/n$, in particular under the hypotheses of
Theorem~\ref{thm:main}(4) or of Theorem~\ref{thm:small}, also
\begin{equation}\label{eq:mixupper}
 t_{\mathrm{mix}}(\varepsilon)\le\Bigl\lceil\frac n2
 \Bigl(\log\frac1{2\varepsilon}+\frac12\log(M-1)\Bigr)\Bigr\rceil .
\end{equation}
\end{corollary}

No sharpness of \eqref{eq:mixlower} or \eqref{eq:mixupper} is asserted.

\paragraph{History and significance.}
Strona, Nappo, Boccacci, Fattorini and San-Miguel-Ayanz~\cite{SNBFS14}
introduced Curveball trades for binary matrices with fixed margins;
Verhelst~\cite{Verhelst08} had developed trade-based sampling for the same
state space. Carstens, Berger and Strona defined undirected
Curveball trades and proved convergence to the uniform
distribution~\cite[Theorem~11]{CBS18}; they left theoretical mixing bounds
open. Carstens and Kleer~\cite{CK18} bounded the relaxation time of the
Curveball chain for binary matrices relative to that of the switch chain;
their comparison gives absolute bounds only where switch-chain bounds are
known, and it does not treat undirected graphs.
Caputo~\cite{Caputo08} studied binary collision processes with a conservation
law, of which row trades on binary matrices with fixed column sums are an
instance, by expanding the square of the generator over triples of particles.
Fu, Qin and Wang~\cite[Lemma~4.1 and Theorem~5.1]{FQW26} used this expansion
to prove that the row-trade chain on binary matrices with arbitrary margins
has gap at least $1/\binom m2$, where $m$ is the number of rows. The
manuscript~\cite{OAI26}, which has not been peer reviewed and whose main
mixing theorem has an agent-written Lean formalization upstream, developed
the pair-fiber operators for undirected graphs.
For the switch chain on regular graphs, Cooper, Dyer and
Greenhill~\cite{CDG07,CDG12} proved polynomial mixing, and Tikhomirov and
Youssef~\cite{TY23} proved sharp Poincar\'e bounds for regular bipartite graphs
in a specified degree range. For every graphical degree sequence with $M>1$,
\cite[Theorem~1.1]{S26} gives $\operatorname{gap}(K)\ge2/\binom n2$, with
equality whenever $M=2$; how much larger the gap is for other
sequences was left open there. Theorem~\ref{thm:main} shows that on regular
graphs the gap is at most $2/n$, that $n-1$ is an eigenvalue of $H$ for every
$n$ and $d$, and that every eigenvalue of $H$ below $n-1$ on $\mathbf1^\perp$,
if one exists, has a rooted-class eigenfunction. Theorem~\ref{thm:small}
excludes such eigenvalues for $n\le\NMAX$ and for most degrees at $n=12$. Whether
$\operatorname{gap}(K)=2/n$ for every regular degree sequence is open
(Section~\ref{op:section}).

\paragraph{Proof strategy.}
On a regular fiber the two endpoints of a trade have the same number of
singleton neighbors, so a trade assigns each singleton to either endpoint with
probability $1/2$. Summing this identity over pairs gives part~(1). The same
balance shows that exchanging the labels of the two endpoints is itself a
trade. Hence every function fixed by $E_a$ is fixed by $T_a$, which gives
$h_a\succeq\tfrac12(I-T_a)$ and part~(2); the class sum of transpositions acts
on $U_\lambda$ as the scalar $\operatorname{cont}(\lambda)$, and an elementary
bound on contents gives the constants. The two components not covered by
part~(2) are generated by the functions fixed by a point stabilizer, which
gives part~(3). For Theorem~\ref{thm:small} we compute the action of $H$ on
rooted-class indicators in rational arithmetic and decide the resulting matrix
inequality by an integer congruence certificate. The comparison with the random transpositions of
vertex labels is the new step; the pair-fiber operators and the kernel of $H$
are taken from~\cite{OAI26,S26}.

\paragraph{Organization.}
Section~\ref{sec:setup} fixes the operators and the relabeling action.
Section~\ref{edg:section} proves the eigenvalue identity.
Section~\ref{cmp:section} proves the comparison with random transpositions and
the bounds on isotypic components. Section~\ref{root:section} proves the
reduction to rooted-class functions, Theorem~\ref{thm:main} and
Corollary~\ref{cor:mixing}. Section~\ref{cert:section} describes the exact
computation behind Theorem~\ref{thm:small}. Section~\ref{exp:section} locates
the slack in the expansion of $H^2$ used in~\cite{S26} and shows that it cannot
be removed term by term. Section~\ref{bin:section} records the corresponding
statements for binary matrices with equal row sums, where the equality
corresponding to $\operatorname{gap}(K)=2/n$ fails for some margins with
equal column sums, and
Section~\ref{op:section} states open problems. All proof ingredients are
supplied in the paper, except the exact computations of
Section~\ref{cert:section}, of Proposition~\ref{exp:obstruction}(2) and of
the binary-matrix examples of Section~\ref{bin:section}, which are carried
out by ancillary programs.

\section{Operators and symmetry}\label{sec:setup}
This section fixes the operator conventions, describes the pair fibers of a
regular graph, and records the two facts about the symmetric group that the
later sections use. All operators act on real $L^2(\pi)$. Write $A\succeq B$
when $A-B$ is self-adjoint and has nonnegative quadratic form, and let
$\Pi f=\pi(f)\mathbf1$.

\paragraph{Pair fibers.}
Fix $G\in\Omega$ and a pair $a=\{x,w\}$. A vertex $y\notin a$ is \emph{common}
if it is adjacent to both $x$ and $w$, and a \emph{singleton} if it is adjacent
to exactly one of them. Let $S_a(G)$ be the set of singletons. The
\emph{$a$-fiber} of $G$ consists of the graphs in $\Omega$ that agree with $G$
on every pair of vertices other than the pairs joining $a$ to $S_a(G)$. These
are the graphs obtainable from $G$ by a Curveball trade at $a$, and $E_a$ is
conditional uniform averaging on the $a$-fibers~\cite[Section~2]{S26}. Hence
$E_a$ and $h_a$ are orthogonal projections, $K$ is symmetric with nonnegative
spectrum, and \eqref{eq:operators} holds. If $x$ has $c$ common neighbors and
$q$ singleton neighbors, then $d=c+q+\mathbf1[xw]$, and the same identity
holds for $w$ with the same $c$. Therefore
\begin{equation}\label{eq:balance}
 |S_a(G)|=2q,\qquad q=d-c-\mathbf1[xw],
\end{equation}
and each endpoint has exactly $q$ singleton neighbors. A graph in the fiber is
determined by the set $X\subseteq S_a(G)$ of singletons adjacent to $x$, which
has $q$ elements, and every $q$-subset occurs: the resulting graph is simple,
agrees with $G$ away from $a$, and gives each vertex its former degree. The
fiber is thus in bijection with the $q$-subsets of a $2q$-set.

\paragraph{Kernel.}
The kernel of $H$ consists of the constants, for every graphical degree
sequence~\cite[Lemma~6.3]{S26}. In particular $\operatorname{gap}(H)>0$ when
$M>1$. The set $\Omega$ is nonempty, since it contains a circulant graph. The
uniform measure is invariant under relabeling of the vertices, and $S_n$ is
transitive on pairs. So every pair
is an edge of a uniform element of $\Omega$ with probability
\begin{equation}\label{eq:density}
 p=\frac{nd/2}{\binom n2}=\frac d{n-1}\in(0,1),
\end{equation}
so $M\ge2$.

\paragraph{Relabeling.}
For $\sigma\in S_n$, the graph $\sigma G$ has the edges $\sigma(i)\sigma(j)$
for $ij\in G$. Relabeling preserves $\Omega$ and $\pi$, so $T_\sigma$ is an
orthogonal representation of $S_n$ on $L^2(\pi)$. Relabeling by $\sigma$ maps
the $a$-fiber of $G$ onto the $\sigma(a)$-fiber of $\sigma G$. Hence
\begin{equation}\label{eq:equivariance}
 T_\sigma E_aT_\sigma^{-1}=E_{\sigma(a)},\qquad T_\sigma H=HT_\sigma .
\end{equation}
The irreducible complex representations of $S_n$ are the Specht modules
$S^\lambda$ indexed by the partitions $\lambda$ of
$n$~\cite[Theorem~2.4.6]{Sagan01}. They are defined over $\mathbb Q$, so the
same modules over $\mathbb R$ are the irreducible real representations. The space $L^2(\pi)$ is the orthogonal
direct sum of its isotypic components $U_\lambda$, and each $U_\lambda$ is
invariant under $H$ by \eqref{eq:equivariance}. We use two facts.
\begin{enumerate}
\item The sum of all transpositions is central in the group algebra and acts
on $S^\lambda$ as the scalar $\operatorname{cont}(\lambda)$; this is
Frobenius's formula for the character at a
transposition~\cite[Lemma~7]{DS81}. Consequently
\begin{equation}\label{eq:content}
 L_T=\Bigl[\binom n2-\operatorname{cont}(\lambda)\Bigr]I
 \quad\text{on }U_\lambda .
\end{equation}
\item The permutation representation of $S_n$ on $\mathbb R^n$ is
$S^{(n)}\oplus S^{(n-1,1)}$. By Frobenius reciprocity, the space of vectors of
$S^\lambda$ fixed by the stabilizer of the point $n$ has dimension one for
$\lambda\in\{(n),(n-1,1)\}$ and dimension zero for every other $\lambda$.
\end{enumerate}

\section{Edge indicators}\label{edg:section}
This section proves that every centered edge indicator is an eigenfunction of
$H$ with eigenvalue $n-1$. It uses only the balance \eqref{eq:balance} of a
regular fiber. Section~\ref{root:section} uses the result for the upper bound
in Theorem~\ref{thm:main}.

\begin{lemma}[Balanced trades]\label{edg:marginal}
Let $a=\{x,w\}$ and $y\notin a$. Then
\[
 E_a\mathbf1[xy]=\tfrac12\bigl(\mathbf1[xy]+\mathbf1[wy]\bigr).
\]
If a pair $e$ equals $a$ or is disjoint from $a$, then
$E_a\mathbf1[e]=\mathbf1[e]$.
\end{lemma}

\begin{proof}
Fix an $a$-fiber. If $y$ is common or adjacent to neither endpoint, then
$\mathbf1[xy]=\mathbf1[wy]$ is constant on the fiber, and both sides agree. If
$y$ is a singleton, then $\mathbf1[xy]+\mathbf1[wy]=1$ on the fiber. By
\eqref{eq:balance}, the set of singletons adjacent to $x$ is a uniform
$q$-subset of the $2q$ singletons, which contains $y$ with probability $1/2$.
So $E_a\mathbf1[xy]=1/2$ on the fiber. The last assertion holds because a
trade at $a$ changes only pairs joining $a$ to a singleton.
\end{proof}

\begin{proposition}[Edge eigenfunctions]\label{edg:eigen}
For every pair $\{x,y\}$,
\[
 H\mathbf1[xy]=(n-1)\mathbf1[xy]-d .
\]
With $p=d/(n-1)$, the function $g_{xy}=\mathbf1[xy]-p$ is nonzero, has mean
zero, and satisfies $Hg_{xy}=(n-1)g_{xy}$.
\end{proposition}

\begin{proof}
By Lemma~\ref{edg:marginal}, $h_a\mathbf1[xy]=0$ unless $a$ meets $\{x,y\}$ in
exactly one vertex. For $w\notin\{x,y\}$ the pairs $\{x,w\}$ and $\{y,w\}$
contribute
\[
 \tfrac12\bigl(\mathbf1[xy]-\mathbf1[wy]\bigr)
 +\tfrac12\bigl(\mathbf1[xy]-\mathbf1[xw]\bigr).
\]
The vertex $y$ has $d-\mathbf1[xy]$ neighbors other than $x$, and $x$ has
$d-\mathbf1[xy]$ neighbors other than $y$. Summing over the $n-2$ choices of
$w$ gives
\[
 H\mathbf1[xy]=(n-2)\mathbf1[xy]-\bigl(d-\mathbf1[xy]\bigr)
 =(n-1)\mathbf1[xy]-d .
\]
Since $H\mathbf1=0$, it follows that $Hg_{xy}=(n-1)\mathbf1[xy]-d=(n-1)g_{xy}$.
By \eqref{eq:density}, $g_{xy}$ has mean zero and takes the two values $1-p$
and $-p$, each with positive probability, so it is nonzero.
\end{proof}

\begin{lemma}[Complement]\label{edg:complement}
Complementation $G\mapsto\overline G$ is a bijection from $\Omega_{n,d}$ onto
$\Omega_{n,n-1-d}$ which maps each $a$-fiber onto an $a$-fiber. Hence the
operators $H$ on the two spaces are unitarily equivalent by a map that
commutes with relabeling, and they have the same eigenvalues on rooted-class
functions.
\end{lemma}

\begin{proof}
Let $a=\{x,w\}$. A vertex $y\notin a$ is a singleton for $G$ exactly when it is
a singleton for $\overline G$, and then it is adjacent in $\overline G$ to the
endpoint it is not adjacent to in $G$. The two graphs have complementary
edges on every pair. So two graphs agree off the pairs joining $a$ to the
singletons exactly when their complements do, and the fibers correspond.
Composition with the bijection is therefore a unitary map
$L^2(\Omega_{n,n-1-d})\to L^2(\Omega_{n,d})$ that intertwines the operators
$E_a$, hence $H$. Complementation commutes with relabeling, so the map
preserves rooted-class functions.
\end{proof}

\section{Comparison with random transpositions}\label{cmp:section}
This section proves $H\succeq\tfrac12L_T$ and the resulting bounds on the
isotypic components. It uses the balance
\eqref{eq:balance} and the scalar \eqref{eq:content}; no estimate on the
Johnson scheme is needed. Section~\ref{root:section} consumes
Proposition~\ref{cmp:sectors}.

\begin{lemma}[The exchange is a trade]\label{cmp:swap}
Let $a=\{x,w\}$ and let $\tau$ be the transposition of $x$ and $w$. For every
$G\in\Omega$, the graph $\tau G$ lies in the $a$-fiber of $G$.
\end{lemma}

\begin{proof}
Relabeling by $\tau$ fixes every pair disjoint from $a$ and the pair $a$
itself. For $y\notin a$ it exchanges the pairs $xy$ and $wy$. So $\tau G$
agrees with $G$ at every common vertex and at every vertex adjacent to
neither endpoint, and the singletons adjacent to $x$ in $\tau G$ are those
adjacent to $w$ in $G$. Thus $\tau G$ agrees with $G$ off the pairs joining
$a$ to $S_a(G)$, and $\tau G\in\Omega$. This is the definition of the
$a$-fiber. In the description of Section~\ref{sec:setup}, $\tau$ replaces the
set $X$ by its complement in $S_a(G)$, which has $q$ elements by
\eqref{eq:balance}.
\end{proof}

\begin{proposition}[Transposition comparison]\label{cmp:main}
For every pair $a$, put $Q_a=\tfrac12(I+T_a)$. Then $Q_a-E_a$ is an orthogonal
projection,
\begin{equation}\label{cmp:local}
 h_a=\tfrac12(I-T_a)+(Q_a-E_a)\succeq\tfrac12(I-T_a),
\end{equation}
and $H=\tfrac12L_T+R$, where $R=\sum_a(Q_a-E_a)\succeq0$ commutes with every
$T_\sigma$.
\end{proposition}

\begin{proof}
The operator $Q_a$ is the orthogonal projection onto the functions fixed by
$T_a$, since $T_a$ is an orthogonal involution. A function in the range of
$E_a$ is constant on every $a$-fiber. By Lemma~\ref{cmp:swap} it takes the
same value at $G$ and $\tau G$, so it is fixed by $T_a$. Hence the range of
$E_a$ is contained in the range of $Q_a$. For orthogonal projections this
gives $Q_aE_a=E_aQ_a=E_a$, so $Q_a-E_a$ is the orthogonal projection onto the
orthogonal complement of the range of $E_a$ inside the range of $Q_a$.
Then $h_a=(I-Q_a)+(Q_a-E_a)$ and $I-Q_a=\tfrac12(I-T_a)$, which is
\eqref{cmp:local}. Summing over $a$ gives $H=\tfrac12L_T+R$. Both $H$ and
$L_T$ commute with every $T_\sigma$, by \eqref{eq:equivariance} and because
conjugation permutes the transpositions; so does $R$.
\end{proof}

\begin{lemma}[Content bound]\label{cmp:content}
Let $\lambda$ be a partition of $n\ge4$ with $\lambda_1\le n-2$. Then
\[
 \operatorname{cont}(\lambda)\le\binom{n-2}2-1,\qquad
 \binom n2-\operatorname{cont}(\lambda)\ge2n-2,
\]
with equality in both for $\lambda=(n-2,2)$.
\end{lemma}

\begin{proof}
Put $\lambda_1=n-k$, so that $2\le k\le n-1$ and the rows $i\ge2$ contain $k$
cells in total. Since $\binom u2+\binom v2\le\binom{u+v}2$ for nonnegative
integers and $i-1\ge1$ in these rows,
\[
 \operatorname{cont}(\lambda)
 =\binom{n-k}2+\sum_{i\ge2}\Bigl[\binom{\lambda_i}2-(i-1)\lambda_i\Bigr]
 \le\binom{n-k}2+\binom k2-k=:\phi(k).
\]
The function $\phi$ is a convex quadratic polynomial in $k$, and
\[
 \phi(2)=\binom{n-2}2-1,\qquad
 \phi(n-1)=\binom{n-1}2-(n-1)=\frac{(n-1)(n-4)}2=\binom{n-2}2-1 .
\]
So $\phi(k)\le\binom{n-2}2-1$ for $2\le k\le n-1$. The second inequality
follows from $\binom n2-\binom{n-2}2=2n-3$. For $\lambda=(n-2,2)$ the contents
of the second row are $-1$ and $0$, so
$\operatorname{cont}(\lambda)=\binom{n-2}2-1$.
\end{proof}

\begin{proposition}[Isotypic components]\label{cmp:sectors}
For every partition $\lambda$ of $n$,
\[
 H\succeq\frac12\Bigl[\binom n2-\operatorname{cont}(\lambda)\Bigr]I
 \quad\text{on }U_\lambda .
\]
The right side is at least $n-1$ when $\lambda\notin\{(n),(n-1,1)\}$, equals
$n/2$ for $\lambda=(n-1,1)$, and equals $0$ for $\lambda=(n)$.
\end{proposition}

\begin{proof}
The first assertion follows from Proposition~\ref{cmp:main} and
\eqref{eq:content}, because $U_\lambda$ is invariant under $H$ and $L_T$. A
partition other than $(n)$ and $(n-1,1)$ has $\lambda_1\le n-2$, and
Lemma~\ref{cmp:content} applies. Finally
$\operatorname{cont}(n-1,1)=\binom{n-1}2-1$ and
$\operatorname{cont}(n)=\binom n2$.
\end{proof}

\begin{remark}[Unequal degrees]\label{cmp:general}
Lemma~\ref{cmp:swap} and \eqref{cmp:local} hold on $\Omega_d$ for an arbitrary
graphical degree sequence $d$ and every pair $a=\{x,w\}$ with $d_x=d_w$, with
the same proofs: equal degrees give \eqref{eq:balance} for that pair, and the
transposition of $x$ and $w$ preserves $\Omega_d$. Pairs of unequal degree
admit no such comparison, since the exchange does not preserve $\Omega_d$.
\end{remark}

\section{Rooted-class functions and the proof of Theorem~\ref{thm:main}}\label{root:section}
This section reduces the two isotypic components left open by
Proposition~\ref{cmp:sectors} to the functions fixed by a point stabilizer,
and assembles Theorem~\ref{thm:main} and Corollary~\ref{cor:mixing}. Let
$S_{n-1}$ denote the stabilizer of the vertex $n$ in $S_n$, and let
\[
 \mathcal V=\{f\in L^2(\pi):T_\sigma f=f\text{ for every }\sigma\in S_{n-1}\}
\]
be the space of rooted-class functions. Its dimension is the number of orbits
of $S_{n-1}$ on $\Omega$, that is, the number of isomorphism classes of
$d$-regular graphs on $n$ vertices with one distinguished vertex.

\begin{proposition}[Rooted reduction]\label{root:reduction}
Put $\mathcal W=U_{(n)}\oplus U_{(n-1,1)}$. Then $\mathcal V\subseteq\mathcal W$,
$H\mathcal V\subseteq\mathcal V$, and $H$ has the same set of eigenvalues on
$\mathcal W\cap\mathbf1^\perp$ as on $\mathcal V\cap\mathbf1^\perp$.
\end{proposition}

\begin{proof}
The orthogonal projection onto $U_\lambda$ commutes with every $T_\sigma$. So
each isotypic component of a function $f\in\mathcal V$ is fixed by $S_{n-1}$.
By the second fact of Section~\ref{sec:setup}, a module isomorphic to a direct
sum of copies of $S^\lambda$ has no nonzero $S_{n-1}$-fixed vector unless
$\lambda\in\{(n),(n-1,1)\}$. Hence $f\in\mathcal W$. The inclusion
$H\mathcal V\subseteq\mathcal V$ follows from \eqref{eq:equivariance}, and the
constants lie in $\mathcal V$, so $H$ preserves
$\mathcal V\cap\mathbf1^\perp$.

Every eigenvalue of $H$ on $\mathcal V\cap\mathbf1^\perp$ is an eigenvalue on
$\mathcal W\cap\mathbf1^\perp$. Conversely, let $\theta$ be an eigenvalue of
$H$ on $\mathcal W\cap\mathbf1^\perp$. Its eigenspace there is a nonzero
subspace invariant under every $T_\sigma$, so it contains a submodule
isomorphic to $S^{(n)}$ or to $S^{(n-1,1)}$. By the same fact this submodule
contains a nonzero vector fixed by $S_{n-1}$, which is an eigenfunction in
$\mathcal V\cap\mathbf1^\perp$ with eigenvalue $\theta$.
\end{proof}

\begin{lemma}[One rooted class]\label{root:single}
If $d\in\{1,n-2\}$, then every rooted-class function is constant.
\end{lemma}

\begin{proof}
Let $d=1$, so that $n$ is even and $\Omega$ is the set of perfect matchings
of $[n]$. Given two perfect matchings, map the partner of $n$ in the first to
the partner of $n$ in the second, and the remaining $n/2-1$ edges of the first
bijectively to those of the second. This defines a permutation fixing $n$
that carries the first matching to the second. So $S_{n-1}$ is transitive on
$\Omega$ and $\mathcal V$ consists of constants. The case $d=n-2$ follows
because complementation commutes with relabeling
(Lemma~\ref{edg:complement}).
\end{proof}

\begin{proof}[Proof of Theorem~\ref{thm:main}]
Part~(1) is Proposition~\ref{edg:eigen}; the bounds \eqref{eq:upper} follow
because $g_{xy}$ is a nonzero element of $\mathbf1^\perp$ and
$\operatorname{gap}(K)=\operatorname{gap}(H)/\binom n2$ with
$(n-1)/\binom n2=2/n$. Part~(2) is Propositions~\ref{cmp:main}
and~\ref{cmp:sectors}.

For part~(3), the space $\mathbf1^\perp$ is the orthogonal direct sum of
$\mathcal W\cap\mathbf1^\perp$ and the components $U_\lambda$ with
$\lambda\notin\{(n),(n-1,1)\}$, and $H$ preserves each summand. On the latter
components every eigenvalue is at least $n-1$ by
Proposition~\ref{cmp:sectors}. On $\mathcal W\cap\mathbf1^\perp$ the least
eigenvalue is $\mu$ by Proposition~\ref{root:reduction}, with $\mu=\infty$
when $\mathcal V\cap\mathbf1^\perp=0$, in which case
$\mathcal W\cap\mathbf1^\perp=0$ as well. Hence
$\operatorname{gap}(H)\ge\min\{n-1,\mu\}$. Conversely $n-1$ is an eigenvalue
of $H$ on $\mathbf1^\perp$ by part~(1), and so is $\mu$ when it is finite. This
proves \eqref{eq:reduction}.

For part~(4), Lemma~\ref{root:single} gives $\mu=\infty$, and
\eqref{eq:reduction} gives $\operatorname{gap}(H)=n-1$.
\end{proof}

\begin{remark}[Perfect matchings]\label{root:matchings}
For $d=1$ every fiber has at most two elements, namely $G$ and $\tau G$ in the
notation of Lemma~\ref{cmp:swap}, so $R=0$ and $H=\tfrac12L_T$. Diaconis and
Holmes determine the spectrum of this walk: the
permutation module on perfect matchings is the sum of the $S^{2\nu}$ over the
partitions $\nu$ of $n/2$, each once~\cite[Theorem~2]{DH02}, so by
\eqref{eq:content} the eigenvalue of $H$ on $S^{2\nu}$ is
$\tfrac12[\binom n2-\operatorname{cont}(2\nu)]$; their Corollary~1 states
the eigenvalues in the normalization of their kernel. The value $n-1$
is attained at $2\nu=(n-2,2)$. In general the functions $g_{xy}$ span a
quotient of the permutation module on pairs,
$S^{(n)}\oplus S^{(n-1,1)}\oplus S^{(n-2,2)}$. The sum of all $g_{xy}$ and the
sums $\sum_{y\ne x}g_{xy}$ vanish, because a $d$-regular graph has $nd/2$
edges and $d$ edges at $x$. These relations span the first two summands, so
every $g_{xy}$ lies in $U_{(n-2,2)}$, where the bound of
Proposition~\ref{cmp:sectors} is attained.
\end{remark}

\begin{proof}[Proof of Corollary~\ref{cor:mixing}]
The relaxation bound restates \eqref{eq:upper}. Let $g=g_{xy}$ and
$\theta=1-2/n$, so that $Kg=\theta g$ by Theorem~\ref{thm:main}(1) and
\eqref{eq:operators}, and $0<\theta<1$. Choose $G$ with
$|g(G)|=\max|g|>0$. Since $\pi(g)=0$,
\[
 \theta^s|g(G)|=|K^sg(G)|
 =\Bigl|\sum_{G'}\bigl(K^s(G,G')-\pi(G')\bigr)g(G')\Bigr|
 \le2d(s)\max|g| .
\]
So $d(s)\ge\theta^s/2$, and $d(s)\le\varepsilon$ forces
$s\log(1/\theta)\ge\log(1/(2\varepsilon))$. Since
$\log(1/\theta)\le(1-\theta)/\theta$ and $\theta/(1-\theta)=n/2-1$, this is
\eqref{eq:mixlower}; the argument is that
of~\cite[Theorem~12.5]{LPW17}.

For \eqref{eq:mixupper}, $K$ is symmetric with nonnegative spectrum, so its
second largest eigenvalue in absolute value is $1-\operatorname{gap}(K)$. The
standard $L^2$ bound for a symmetric kernel with uniform stationary law gives
$d(s)\le\tfrac12\sqrt{M-1}\,(1-\operatorname{gap}(K))^s
\le\tfrac12\sqrt{M-1}\,e^{-2s/n}$, as in the proof
of~\cite[Corollary~1.2]{S26}. The right side is at most $\varepsilon$ when
$s\ge\tfrac n2[\log(1/(2\varepsilon))+\tfrac12\log(M-1)]$.
\end{proof}

\section{Exact computation for small cases}\label{cert:section}
This section reduces Theorem~\ref{thm:small} to the positive definiteness of
one rational matrix for each pair $(n,d)$ and reports the computation. By
Theorem~\ref{thm:main}(3) it suffices to show $\mu>n-1$, and by
Lemma~\ref{edg:complement} it suffices to treat $d\le(n-1)/2$. The
computation is exact; no floating-point quantity enters the proof.

Let $O_1,\ldots,O_r$ be the orbits of $S_{n-1}$ on $\Omega$, with sizes
$w_1,\ldots,w_r$ and representatives $G_1,\ldots,G_r$. Their indicators form
a basis of $\mathcal V$. Write $F_a(G)$ for the $a$-fiber of $G$ and define
the $r\times r$ rational matrix
\begin{equation}\label{cert:matrix}
 A_{cc'}=\sum_{a\in\binom{[n]}2}
 \Bigl(\mathbf1[c=c']-\frac{|F_a(G_c)\cap O_{c'}|}{|F_a(G_c)|}\Bigr).
\end{equation}

\begin{lemma}[Quotient criterion]\label{cert:criterion}
For $f=\sum_cf_c\mathbf1_{O_c}\in\mathcal V$, the function $Hf$ takes the
value $\sum_{c'}A_{cc'}f_{c'}$ on $O_c$. With $W=\operatorname{diag}(w)$ and
$w=(w_1,\ldots,w_r)^{\mathsf T}$, the matrix
\begin{equation}\label{cert:B}
 B=W\bigl(A-(n-1)I\bigr)+\frac nM\,ww^{\mathsf T}
\end{equation}
is symmetric, and it is positive definite if and only if $\mu>n-1$.
\end{lemma}

\begin{proof}
By definition $Hf(G_c)=\sum_a[f(G_c)-|F_a(G_c)|^{-1}\sum_{G'\in F_a(G_c)}f(G')]$,
which is the stated value, and $Hf\in\mathcal V$ is constant on $O_c$ by
Proposition~\ref{root:reduction}. For $f,g\in\mathcal V$ with coordinate
vectors $f,g\in\mathbb R^r$,
\[
 g^{\mathsf T}WAf=M\langle g,Hf\rangle,\qquad
 g^{\mathsf T}Wf=M\langle g,f\rangle,\qquad w^{\mathsf T}f=M\pi(f).
\]
The first form is symmetric because $H$ is self-adjoint, so $B$ is
symmetric. Write $f=f_0+\pi(f)\mathbf1$ with $\pi(f_0)=0$. Since
$H\mathbf1=0$,
\[
 f^{\mathsf T}Bf
 =M\bigl[\langle f_0,Hf_0\rangle-(n-1)\langle f_0,f_0\rangle+\pi(f)^2\bigr].
\]
This is positive for every $f\ne0$ exactly when
$\langle f_0,Hf_0\rangle>(n-1)\langle f_0,f_0\rangle$ for every nonzero
centered $f_0\in\mathcal V$, that is, when $\mu>n-1$.
\end{proof}

\begin{proof}[Proof of Theorem~\ref{thm:small}]
The upper bound is Theorem~\ref{thm:main}(1). For $d\in\{1,n-2\}$ the lower
bound is Theorem~\ref{thm:main}(4). For $n=5$ and $d=2$ the only graphs are
the $5$-cycles, which form one rooted class, so $\mu=\infty$ and
Theorem~\ref{thm:main}(3) gives the lower bound.
For the remaining pairs with $d\le(n-1)/2$, listed in
Table~\ref{cert:table}, an ancillary program computes the orbits, the
matrix \eqref{cert:matrix} and the sizes $w_c$ in integer and rational
arithmetic. It then exhibits an integer matrix $Q$ such that
$Q^{\mathsf T}B'Q$ has positive diagonal entries and is strictly diagonally
dominant, where $B'$ is a positive integer multiple of $B$ with integer
entries; all entries of the product are computed in integer arithmetic. A
symmetric strictly diagonally dominant matrix with positive diagonal is
positive definite by Gershgorin's theorem. Then $Q$ is nonsingular, since
$Qv=0$ with $v\ne0$ would give $v^{\mathsf T}Q^{\mathsf T}B'Qv=0$, and $B$
is congruent to a positive definite matrix, hence positive definite.
Lemma~\ref{cert:criterion} gives $\mu>n-1$, and
Theorem~\ref{thm:main}(3) gives $\operatorname{gap}(H)=n-1$. The pairs with
$d>(n-1)/2$ follow from Lemma~\ref{edg:complement}. The matrix $Q$ is an integer diagonal matrix times a rounding of a
floating-point eigenvector matrix of a diagonal rescaling of $B'$; this
affects
only whether a certificate is found, and the verification of the certificate
uses no floating-point quantity.
\end{proof}

The program finds the orbits by a search from one graph along Curveball
trades, which reaches every graph because $\ker H=\mathbb R\mathbf1$, and
identifies orbits by a canonical form of the rooted graph. It computes
$w_c=(n-1)!/|\operatorname{Aut}(G_c,n)|$, where $\operatorname{Aut}(G_c,n)$ is
the group of automorphisms of $G_c$ fixing $n$. As internal checks it
verifies the detailed-balance identities $w_cA_{cc'}=w_{c'}A_{c'c}$ exactly
and compares $\sum_cw_c$ with the number $M$ of labeled $d$-regular graphs
computed by an independent recursion on degree sequences. For at most $64$
rooted classes it also decides positive definiteness by fraction-free
Gaussian elimination, with the same result.
Table~\ref{cert:table} lists the sizes and, for orientation, a floating-point
value of $\mu$.

\begin{table}[ht]
\centering
\begin{tabular}{rrrrrr}
\toprule
$n$ & $d$ & $M$ & classes & rooted classes & $\mu$\\
\midrule
$6$ & $2$ & $70$ & $2$ & $2$ & $7.0000$\\
$7$ & $2$ & $465$ & $2$ & $3$ & $9.1667$\\
$8$ & $2$ & $3{,}507$ & $3$ & $4$ & $10.4185$\\
$8$ & $3$ & $19{,}355$ & $6$ & $11$ & $8.2221$\\
$9$ & $2$ & $30{,}016$ & $4$ & $6$ & $11.5237$\\
$9$ & $4$ & $1{,}024{,}380$ & $16$ & $50$ & $9.7832$\\
$10$ & $2$ & $286{,}884$ & $5$ & $8$ & $14.1428$\\
$10$ & $3$ & $11{,}180{,}820$ & $21$ & $68$ & $12.7077$\\
$10$ & $4$ & $66{,}462{,}606$ & $60$ & $286$ & $11.4954$\\
$11$ & $2$ & $3{,}026{,}655$ & $6$ & $11$ & $16.7409$\\
$11$ & $4$ & $5{,}188{,}453{,}830$ & $266$ & $1{,}918$ & $13.3403$\\
$12$ & $2$ & $34{,}944{,}085$ & $9$ & $15$ & $18.4146$\\
$12$ & $3$ & $11{,}555{,}272{,}575$ & $94$ & $510$ & $16.4066$\\
$12$ & $4$ & $480{,}413{,}921{,}130$ & $1{,}547$ & $14{,}577$ & $15.1666$\\
\bottomrule
\end{tabular}
\caption{The cases of Theorem~\ref{thm:small} with $2\le d\le(n-1)/2$ and
more than one rooted class; the pair $n=12$, $d=5$ is not covered. Classes are isomorphism classes of $d$-regular
graphs on $n$ vertices; rooted classes are the orbits of $S_{n-1}$. In every
row the matrix $B$ of \eqref{cert:B} is positive definite by the exact
certificate described above. The last column is a floating-point value of $\mu$, rounded to
four decimals; it is not used in the proof.}\label{cert:table}
\end{table}

\section{The expansion of \texorpdfstring{$H^2$}{H squared}}\label{exp:section}
The universal bound $H\succeq2(I-\Pi)$ of~\cite{S26} comes from an expansion
of $H^2$ into a triple term and a disjoint-pair term, each bounded
separately. The expansion is due to
Caputo~\cite[Equations~(1.13) to~(1.15)]{Caputo08} and was applied to pair
fibers of graphs in~\cite[Proposition~7.1]{OAI26}. This section computes both terms on the edge indicators of a
regular fiber, where $S$ and $D$ both equal $\tfrac{n-2}2H$, against the
universal bounds $S\succeq0$ and $D\succeq H$, and shows that
on regular graphs the two terms cannot be bounded separately with constants
that add up to the conjectured gap. Nothing here is used in the proofs of
Theorems~\ref{thm:main} and~\ref{thm:small}.

For a vertex triple $T$ let $H_T$ be the sum of $h_a$ over the three pairs
$a\subset T$, and let $D$ be the sum of the products $h_ah_b$ over ordered
pairs $(a,b)$ of disjoint vertex pairs. Put $S=\sum_T(H_T^2-H_T)$. Since
$\sum_TH_T=(n-2)H$, the expansion~\cite[Equation~(6.1)]{S26} reads
\begin{equation}\label{exp:expansion}
 H^2=H+S+D .
\end{equation}
For every graphical degree sequence, $S\succeq0$ and
$D\succeq H$~\cite[Lemma~A.4 and Proposition~6.1]{S26}, which gives
$H^2\succeq2H$. Because $H\succeq0$ has kernel $\mathbb R\mathbf1$, the
equality $\operatorname{gap}(H)=n-1$ on a regular fiber is equivalent to
$H^2\succeq(n-1)H$, that is, to
\begin{equation}\label{exp:target}
 S+D\succeq(n-2)H .
\end{equation}

\begin{proposition}[Edge functions]\label{exp:edges}
Let $\mathcal E$ be the span of $\mathbf1$ and the edge indicators
$\mathbf1[e]$. Each $h_a$ preserves $\mathcal E$ and equals
$\tfrac12(I-T_a)$ there. On $\mathcal E$,
\[
 H=\tfrac12L_T,\qquad H_T^2=\tfrac32H_T\ \text{ for every triple }T,\qquad
 S=\frac{n-2}2H,
\]
and $Dg_{xy}=\tfrac{n-2}2Hg_{xy}=\tfrac{(n-1)(n-2)}2g_{xy}$ for every pair
$\{x,y\}$.
\end{proposition}

\begin{proof}
For $a=\{x,w\}$, the relabeling $T_a$ exchanges $\mathbf1[xy]$ and
$\mathbf1[wy]$ for $y\notin a$ and fixes $\mathbf1[e]$ when $e=a$ or
$e\cap a=\varnothing$. Lemma~\ref{edg:marginal} therefore gives
$E_a=\tfrac12(I+T_a)$ on $\mathcal E$, which proves the first assertion and
$H=\tfrac12L_T$. On $\mathcal E$ we then have
$H_T=\tfrac12(3I-C_T)$, where $C_T$ is the sum of the three transpositions of
the vertices of $T$. Under the group of permutations of $T$, the space $\mathcal E$ is a quotient
of the sum of a trivial representation and the permutation representation on
pairs. The orbits on pairs have size one or three, and
an orbit of size three carries the permutation representation on three
points, which is the sum of the trivial and the two-dimensional irreducible
representation. On these $C_T$ acts as $3$ and as $0$, and the sign
representation does not occur. So $H_T$ has eigenvalues $0$ and $3/2$ on
$\mathcal E$, and $H_T^2=\tfrac32H_T$. Summing $H_T^2-H_T=\tfrac12H_T$ over
$T$ gives $S=\tfrac{n-2}2H$. Finally \eqref{exp:expansion} and
Proposition~\ref{edg:eigen} give
$Dg_{xy}=[(n-1)^2-(n-1)-\tfrac{(n-1)(n-2)}2]g_{xy}$.
\end{proof}

At the functions $g_{xy}$ the two sides of \eqref{exp:target} are equal, and
$S$ and $D$ contribute equally. The next proposition shows that neither
$S\succeq\tfrac{n-2}2H$ nor any other pair of separate bounds summing to
\eqref{exp:target} holds in general.

\begin{proposition}[Obstruction to separate bounds]\label{exp:obstruction}
\begin{enumerate}
\item On $\Omega_{7,2}$ and $T=\{1,2,3\}$, the operator $H_T$ has the
eigenvalue $1$. So the inequality $H_T^2\succeq H_T$
of~\cite[Lemma~A.4]{S26} is attained on a regular fiber, and
$H_T^2\succeq\tfrac32H_T$ fails.
\item On $\Omega_{8,3}$ the number $t$ of triangles of a graph satisfies
\[
 \frac{\langle t,St\rangle}{\langle t,Ht\rangle}=\frac{2314}{835}<3 .
\]
Consequently there are no constants $\alpha,\beta$ with $\alpha+\beta=6$,
$S\succeq\alpha H$ and $D\succeq\beta H$ on $\Omega_{8,3}$.
\end{enumerate}
\end{proposition}

\begin{proof}
For~(1), let $G$ be the $7$-cycle with the cyclic vertex order
$(1,4,2,5,6,3,7)$. No pair inside $T$ is an edge, the vertices $4$ and $7$
each have two neighbors in $T$, and the vertices $5$ and $6$ each have one.
Trades at the pairs inside $T$ preserve these properties and the edge
$\{5,6\}$. The set $\mathcal C$ of graphs reachable from $G$ by such trades
is the set of all assignments of neighborhoods in $T$ to the vertices
$4,5,6,7$ of these sizes in which every vertex of $T$ has degree two; it has
$15$ elements, and functions supported on $\mathcal C$ form an invariant
subspace of $H_T$. Let $c_0$ be the set of graphs in $\mathcal C$ in which
$5$ and $6$ have the same neighbor in $T$, let $c_1$ be the set in which the
neighbor of $5$ in $T$ is adjacent to $4$ and that of $6$ is not, and let
$c_2$ be the set in which the neighbor of $5$ in $T$ is adjacent to $7$ and
that of $6$ is not. These sets have $3$, $6$ and $6$ elements and partition
$\mathcal C$, and they are invariant under the permutations of $T$.
Enumerating the fibers, which have six or two elements, shows that $H_T$
acts on the column vector of coefficients of
$(\mathbf1_{c_0},\mathbf1_{c_1},\mathbf1_{c_2})$ by
\[
 \begin{pmatrix}4/3&-2/3&-2/3\\-1/3&2/3&-1/3\\-1/3&-1/3&2/3\end{pmatrix},
\]
so $f=\mathbf1_{c_1}-\mathbf1_{c_2}$ satisfies $H_Tf=f$. An ancillary program
confirms this identity on $\mathcal C$ in rational arithmetic.
Statement~(2) is a finite computation in rational arithmetic, carried out
by an ancillary program, which evaluates both quadratic forms over the $19{,}355$
elements of $\Omega_{8,3}$; constants do not contribute, since
$H\mathbf1=H_T\mathbf1=0$. The displayed ratio forces $\alpha\le2314/835<3$,
and Proposition~\ref{exp:edges} with $n=8$ forces $\beta\le3$.
\end{proof}

\section{Binary matrices with equal row sums}\label{bin:section}
Curveball trades were introduced for binary matrices with fixed
margins~\cite{SNBFS14}. This section records the statements of
Sections~\ref{edg:section} to~\ref{root:section} in that setting, where the
proofs are the same and the conclusion is different: the comparison leaves
only the functions invariant under row permutations, and on these the bound
can fail when the column sums are unequal. Nothing in this section is used
elsewhere.

Let $m\ge2$ and $k\ge1$, let $1\le r\le k-1$, and let
$c=(c_1,\ldots,c_k)$ be integers with $0<c_j<m$ and $\sum_jc_j=mr$. Let
$\Omega_{r,c}$ be the set of $m\times k$ binary matrices with every row sum
equal to $r$ and column sums $c$, assumed nonempty, with the uniform
measure. For a pair $a=\{x,w\}$ of rows, the $a$-fiber of a matrix $Z$
consists of the matrices in $\Omega_{r,c}$ that agree with $Z$ outside rows
$x$ and $w$. A column is a \emph{singleton} if rows $x$ and $w$ differ in it;
since the two rows have the same sum, there are $2q$ singletons, $q$ with a
one in each row, and the fiber is in bijection with the $q$-subsets of the
singletons. Let $E_a$ average over the $a$-fibers, and put $h_a=I-E_a$,
$H=\sum_ah_a$ over the $\binom m2$ row pairs, and $K=I-H/\binom m2$. Thus
$K$ performs a uniform Curveball trade at a uniformly chosen pair of rows.
The chain is irreducible by Ryser's interchange
theorem~\cite[Theorem~3.1]{Ryser57}, since an interchange is a trade. The
group $S_m$ acts by permuting rows, $L_T$ denotes the sum of $I-T_a$ over the
row transpositions, and a function is \emph{row-symmetric} if it is invariant
under every row permutation.

\begin{proposition}[Row trades]\label{bin:main}
In the setting above:
\begin{enumerate}
\item for every row $x$ and column $j$, the function $\mathbf1[Z_{xj}=1]-c_j/m$
is nonzero and is an eigenfunction of $H$ with eigenvalue $m/2$;
\item $H\succeq\tfrac12L_T$, and $H\succeq\tfrac m2I$ on the orthogonal
complement of the row-symmetric functions;
\item $\operatorname{gap}(H)=\min\{m/2,\mu_{\mathrm{sym}}\}$, where
$\mu_{\mathrm{sym}}$ is the least eigenvalue of $H$ on the centered
row-symmetric functions, and $\mu_{\mathrm{sym}}=\infty$ when there are none.
\end{enumerate}
\end{proposition}

\begin{proof}
(1) As in Lemma~\ref{edg:marginal}, a trade at $\{x,w\}$ gives
$E_a\mathbf1[Z_{xj}=1]=\tfrac12(\mathbf1[Z_{xj}=1]+\mathbf1[Z_{wj}=1])$, and
pairs of rows not containing $x$ fix $\mathbf1[Z_{xj}=1]$. Summing over the
$m-1$ rows $w\ne x$ gives
\[
 H\mathbf1[Z_{xj}=1]=\tfrac{m-1}2\mathbf1[Z_{xj}=1]
 -\tfrac12\bigl(c_j-\mathbf1[Z_{xj}=1]\bigr)
 =\tfrac m2\mathbf1[Z_{xj}=1]-\tfrac{c_j}2 .
\]
By row symmetry the mean of $\mathbf1[Z_{xj}=1]$ is $c_j/m\in(0,1)$, so the
centered function is nonzero and the claim follows from $H\mathbf1=0$.

(2) Exchanging rows $x$ and $w$ preserves $\Omega_{r,c}$ and replaces the set
of singletons with a one in row $x$ by its complement, so it is a trade. The
proof of Proposition~\ref{cmp:main} applies verbatim and gives
$H\succeq\tfrac12L_T$. On the $\lambda$-isotypic component for a partition
$\lambda\ne(m)$ of $m$, with $\lambda_1=m-k'$ and $1\le k'\le m-1$, the
argument of Lemma~\ref{cmp:content} gives
$\operatorname{cont}(\lambda)\le\binom{m-k'}2+\binom{k'}2-k'$, a convex
function of $k'$ whose values at $k'=1$ and $k'=m-1$ are
$\binom{m-1}2-1$ and $\binom{m-1}2-(m-1)$. Hence
$\binom m2-\operatorname{cont}(\lambda)\ge\binom m2-\binom{m-1}2+1=m$, and
$H\succeq\tfrac m2I$ on every component other than $U_{(m)}$, which is the
space of row-symmetric functions.

(3) The space $\mathbf1^\perp$ is the orthogonal sum of the centered
row-symmetric functions and the components $U_\lambda$ with
$\lambda\ne(m)$, each invariant under $H$. On the latter the least eigenvalue
is at least $m/2$ by~(2), and $m/2$ is an eigenvalue on $\mathbf1^\perp$
by~(1). So the least eigenvalue of $H$ on $\mathbf1^\perp$ is the smaller of
$m/2$ and $\mu_{\mathrm{sym}}$.
\end{proof}

For the margins $m=3$, $r=2$, $c=(1,1,2,2)$ the space $\Omega_{r,c}$ has $15$
elements, and $H$ has the eigenvalue $1<3/2$ with a row-symmetric
eigenfunction; this is the computation in the proof of
Proposition~\ref{exp:obstruction}(1), in which the three rows are the
neighborhoods of the vertices of $T$ outside $T$. So equal row sums alone do
not give $\operatorname{gap}(H)=m/2$.

Equal column sums do not give it either. For $m=k=4$ and all row and column
sums equal to $2$, the space has $90$ elements and
$\mu_{\mathrm{sym}}=4/3<2$. For square matrices with all row and column sums
equal to $2$, an exact computation of the same kind as in
Section~\ref{cert:section}, on the classes of matrices under row
permutations, gives $\mu_{\mathrm{sym}}<m/2$ for $m\in\{4,5,6\}$ and
$\mu_{\mathrm{sym}}>m/2$ for $m\in\{7,8\}$; for $m=k=6$ and all sums equal to
$3$ it gives $\mu_{\mathrm{sym}}<3$. The statement corresponding to
Conjecture~\ref{op:conjecture} is therefore false for binary matrices with
equal row sums and equal column sums. In the graph setting rows and columns are permuted by the same relabelings,
so the space $U_{(n)}$ of invariant functions is smaller than the space of
row-symmetric functions of an adjacency matrix.

The expansion of Section~\ref{exp:section} gives a sufficient condition.
Disjoint row pairs resample disjoint rows, so their projections commute and
$D\succeq0$. Let $m\ge3$, and let $\kappa$ be the least eigenvalue, over all triples $T$
of rows, of $H_T$ on the orthogonal complement of its kernel. Then
$H_T^2\succeq\kappa H_T$, and the expansion
$H^2=\sum_TH_T^2-(m-3)H+D$ with $\sum_TH_T=(m-2)H$ gives
\begin{equation}\label{bin:caputo}
 \operatorname{gap}(H)\ge(\kappa-1)(m-2)+1 ,
\end{equation}
which is Caputo's three-particle bound~\cite[Theorem~3.1]{Caputo08} in the
present normalization. Fu, Qin and Wang~\cite[Theorem~5.1]{FQW26} prove
$\kappa\ge1$ for all margins, which gives $\operatorname{gap}(H)\ge1$; the
example with $m=3$ above attains it. If $\kappa\ge3/2$, then
\eqref{bin:caputo} and Proposition~\ref{bin:main}(1) give
$\operatorname{gap}(H)=m/2$. This condition is not necessary: for $m=k=7$
and all sums equal to $2$, fixing the rows $\{3,4\}$, $\{5,6\}$, $\{5,7\}$,
$\{6,7\}$, written as their sets of ones, leaves three rows with the column
sums $(2,2,1,1,0,0,0)$, so $\kappa=1$, while $\mu_{\mathrm{sym}}>7/2$.

\section{Open problems}\label{op:section}
The equality $\operatorname{gap}(K)=2/n$ is proved here for $d\in\{1,n-2\}$
and in the cases of Theorem~\ref{thm:small}. It is not proved for any other
regular degree sequence. For $n=12$ and $d\in\{5,6\}$ there are $7{,}849$
classes and $83{,}858$ rooted classes; a floating-point computation gives
$\mu\approx14.62>11$, which is evidence and not a proof.

\begin{conjecture}[Regular Curveball gap]\label{op:conjecture}
For every $n\ge4$ and $1\le d\le n-2$ with $nd$ even,
$\operatorname{gap}(H)=n-1$ and $\operatorname{gap}(K)=2/n$ on $\Omega_{n,d}$.
\end{conjecture}

By Theorem~\ref{thm:main}(3), Conjecture~\ref{op:conjecture} is equivalent to
$\mu\ge n-1$, an inequality for functions of the rooted isomorphism class. On
the centered functions of the unrooted isomorphism class, which form
$U_{(n)}\cap\mathbf1^\perp$, the comparison of Proposition~\ref{cmp:main}
gives no information, and the only bound proved for general $n$ and $d$ is
$H\succeq2I$~\cite[Theorem~1.1]{S26}. On $U_{(n-1,1)}$ the proved bound is
$n/2$. A lower bound for $\mu$ of order $n$ would already show that the
relaxation time of $K$ on regular graphs has order $n$. By
\eqref{exp:target} the conjecture is also equivalent to
$S+D\succeq(n-2)H$, and Proposition~\ref{exp:obstruction} shows that this
cannot be obtained from separate bounds on $S$ and $D$. The conjecture has the
form of the statement proved by Caputo, Liggett and
Richthammer~\cite{CLR10} for the interchange process: the gap of the process
is attained on a small explicit invariant subspace, here the span of the
edge indicators. We do not know whether their method applies; conditioning
on the neighborhood of one vertex leaves the class of regular graphs. A closer
analog is the Kac walk and its relatives, where the gap is attained at an
explicit function and is obtained from the three-particle gap by the
expansion of Section~\ref{exp:section}~\cite{Caputo08}; for regular graphs
Proposition~\ref{exp:obstruction}(1) shows that the three-vertex operators do
not have the gap this route requires.

In Table~\ref{cert:table} the computed values of $\mu$ exceed $n-1$ by at
least $1.2$, and the least ratio $\mu/(n-1)$ is about $1.17$, at $n=8$,
$d=3$. We do not know whether $\mu-(n-1)$ is bounded below by a positive
constant, or grows with $n$, over all regular degree sequences.

For a degree sequence with unequal degrees, Remark~\ref{cmp:general} compares
only the trades at pairs of equal degree with transpositions. No analog of
Theorem~\ref{thm:main}(1) is known there, and the gap of $H$ can equal
$2$~\cite[Theorem~1.1]{S26}. For binary matrices with equal row sums,
Proposition~\ref{bin:main} leaves the row-symmetric functions. The
margins for which $\mu_{\mathrm{sym}}\ge m/2$ are not characterized; the
criterion \eqref{bin:caputo} is sufficient and is not necessary.

The bound \eqref{eq:mixupper} loses the factor $\log M$ of the passage from
the spectral gap to total variation. For $d=1$ the walk is a quotient of the
random-transposition walk on $S_n$ and mixes in order $n\log n$
steps~\cite{DS81,DH02}, while \eqref{eq:mixupper} has order $n^2\log n$. An
entropy or log-Sobolev inequality for $K$ on regular graphs is not proved
here.

\section*{Acknowledgments}
The manuscript, its exact arithmetic verification and the computations of
Sections~\ref{cert:section}, \ref{exp:section} and~\ref{bin:section} were produced
with AI assistance. The verification artifacts, the ancillary programs and
the manuscript sources are available in the repository and the archived
record
\begin{center}
\url{https://github.com/sevenevesai/RegularCurveballGap},\\
\url{https://doi.org/10.5281/zenodo.24180369}.
\end{center}

\begin{thebibliography}{99}

\bibitem{Caputo08}
P.~Caputo,
\emph{On the spectral gap of the Kac walk and other binary collision processes},
ALEA. Latin American Journal of Probability and Mathematical Statistics
\textbf{4} (2008), 205--222.
\href{https://alea.impa.br/articles/v4/04-10.pdf}{Primary article}.

\bibitem{CLR10}
P.~Caputo, T.~M.~Liggett, and T.~Richthammer,
\emph{Proof of Aldous' spectral gap conjecture},
Journal of the American Mathematical Society \textbf{23} (2010), no.~3, 831--851.
\href{https://doi.org/10.1090/S0894-0347-10-00659-4}{doi:10.1090/S0894-0347-10-00659-4}.

\bibitem{CBS18}
C.~J.~Carstens, A.~Berger, and G.~Strona,
\emph{A unifying framework for fast randomization of ecological networks with fixed (node) degrees},
MethodsX \textbf{5} (2018), 773--780.\newline
\href{https://doi.org/10.1016/j.mex.2018.06.018}{doi:10.1016/j.mex.2018.06.018}.
Extended version: \href{https://arxiv.org/abs/1609.05137v3}{arXiv:1609.05137v3};
the theorem cited here is numbered as in the extended version.

\bibitem{CK18}
C.~J.~Carstens and P.~Kleer,
\emph{Speeding up Switch Markov Chains for Sampling Bipartite Graphs
with Given Degree Sequence},
Approximation, Randomization, and Combinatorial Optimization.
Algorithms and Techniques (APPROX/RANDOM 2018),
Leibniz International Proceedings in Informatics \textbf{116} (2018), 36:1--36:18.
\href{https://doi.org/10.4230/LIPIcs.APPROX-RANDOM.2018.36}{doi:10.4230/LIPIcs.APPROX-RANDOM.2018.36}.

\bibitem{CDG07}
C.~Cooper, M.~Dyer, and C.~Greenhill,
\emph{Sampling regular graphs and a peer-to-peer network},
Combinatorics, Probability and Computing \textbf{16} (2007), no.~4, 557--593.\newline
\href{https://doi.org/10.1017/S0963548306007978}{doi:10.1017/S0963548306007978}.

\bibitem{CDG12}
C.~Cooper, M.~Dyer, and C.~Greenhill,
\emph{Corrigendum: Sampling regular graphs and a peer-to-peer network},
preprint (2012),
\href{https://arxiv.org/abs/1203.6111v1}{arXiv:1203.6111v1}.

\bibitem{DH02}
P.~Diaconis and S.~P.~Holmes,
\emph{Random walks on trees and matchings},
Electronic Journal of Probability \textbf{7} (2002), paper no.~6, 1--17.
\href{https://doi.org/10.1214/EJP.v7-105}{doi:10.1214/EJP.v7-105}.

\bibitem{DS81}
P.~Diaconis and M.~Shahshahani,
\emph{Generating a random permutation with random transpositions},
Zeitschrift f\"ur Wahrscheinlichkeitstheorie und verwandte Gebiete
\textbf{57} (1981), no.~2, 159--179.
\href{https://doi.org/10.1007/BF00535487}{doi:10.1007/BF00535487}.

\bibitem{FQW26}
W.~Fu, Q.~Qin, and G.~Wang,
\emph{Spectral Gap for the Binary Fixed-Margin Swap Chain},
preprint (2026), \href{https://arxiv.org/abs/2606.22636v2}{arXiv:2606.22636v2}.

\bibitem{LPW17}
D.~A.~Levin and Y.~Peres,
\emph{Markov chains and mixing times}, second edition,
with contributions by E.~L.~Wilmer,
American Mathematical Society, Providence, RI, 2017.
\href{https://doi.org/10.1090/mbk/107}{doi:10.1090/mbk/107}.

\bibitem{OAI26}
OpenAI,
\emph{Polynomial mixing of the switch chain for every graphical degree sequence},
OpenAI Math Release preprint
\href{https://github.com/openai/math/tree/adc7f1241b42e322a6451854ab7e4b4c146bf78a/preprints/Polynomial-Mixing-of-the-Switch-Chain-for-Every-Graphical-Degree-Sequence-September-25-2026/build}{OAI:Polynomial-Mixing-of-the-Switch-Chain-for-Every-Graphical-Degree-Sequence-September-25-2026},
2026.

\bibitem{Ryser57}
H.~J.~Ryser,
\emph{Combinatorial properties of matrices of zeros and ones},
Canadian Journal of Mathematics \textbf{9} (1957), 371--377.
\href{https://doi.org/10.4153/CJM-1957-044-3}{doi:10.4153/CJM-1957-044-3}.

\bibitem{Sagan01}
B.~E.~Sagan,
\emph{The symmetric group. Representations, combinatorial algorithms, and
symmetric functions}, second edition,
Graduate Texts in Mathematics \textbf{203}, Springer, New York, 2001.
\href{https://doi.org/10.1007/978-1-4757-6804-6}{doi:10.1007/978-1-4757-6804-6}.

\bibitem{S26}
M.~T.~Short,
\emph{Sharp spectral gaps for the switch and Curveball chains on every
graphical degree sequence},
preprint (2026).
\href{https://doi.org/10.5281/zenodo.23295359}{doi:10.5281/zenodo.23295359}.

\bibitem{SNBFS14}
G.~Strona, D.~Nappo, F.~Boccacci, S.~Fattorini, and J.~San-Miguel-Ayanz,
\emph{A fast and unbiased procedure to randomize ecological binary matrices
with fixed row and column totals},
Nature Communications \textbf{5} (2014), 4114.
\href{https://doi.org/10.1038/ncomms5114}{doi:10.1038/ncomms5114}.

\bibitem{TY23}
K.~Tikhomirov and P.~Youssef,
\emph{Sharp Poincar\'e and log-Sobolev inequalities for the switch chain on
regular bipartite graphs},
Probability Theory and Related Fields \textbf{185} (2023), 89--184.
\href{https://doi.org/10.1007/s00440-022-01172-7}{doi:10.1007/s00440-022-01172-7}.

\bibitem{Verhelst08}
N.~D.~Verhelst,
\emph{An efficient MCMC algorithm to sample binary matrices with fixed marginals},
Psychometrika \textbf{73} (2008), no.~4, 705--728.
\href{https://doi.org/10.1007/s11336-008-9062-3}{doi:10.1007/s11336-008-9062-3}.

\end{thebibliography}

\end{document}
